% =====================================================================
\section{Construction et déploiement}

\subsection{De la source au WAR}

\begin{lstlisting}[style=shell]
./mvnw clean package        # Windows : mvnw.cmd clean package
\end{lstlisting}

Maven compile les sources avec le \textbf{JDK 21}, copie les ressources
(\texttt{persistence.xml}) et les pages web, puis produit
\texttt{target/gestion-agences.war}. Le \emph{wrapper} (\texttt{mvnw}) est
versionné avec le projet : il télécharge lui-même la bonne version de Maven,
donc la construction ne dépend pas de ce qui est installé sur le poste.

\subsection{Ce que contient le WAR}

\noindent\begin{tabularx}{\textwidth}{|>{\ttfamily}L{5.8cm}|X|}
\hline
\textmd{\textbf{Chemin dans l'archive}} & \textbf{Contenu}\\
\hline
index.jsp, forgot-password.jsp & pages accessibles sans authentification\\
\hline
css/style.css & feuille de style\\
\hline
WEB-INF/*.jsp & vues protégées, inaccessibles par URL directe\\
\hline
WEB-INF/web.xml & fichier de déploiement (page d'accueil)\\
\hline
WEB-INF/jboss-web.xml & racine de contexte \texttt{/gestion-agences}\\
\hline
WEB-INF/classes/ & classes compilées et \texttt{META-INF/persistence.xml}\\
\hline
\end{tabularx}

\vspace{0.5em}
Le fichier \texttt{web.xml} est volontairement minimal : les servlets et le
filtre sont déclarés par \textbf{annotations}
(\texttt{@WebServlet}, \texttt{@WebFilter}), pas dans le descripteur.

\begin{lstlisting}[style=xml]
<jboss-web>
    <context-root>/gestion-agences</context-root>
</jboss-web>
\end{lstlisting}

C'est ce fichier qui explique que l'application réponde sur
\texttt{http://localhost:8080/gestion-agences/} et non sous le nom de
l'archive.

\subsection{Déployer sur WildFly}

Copier le WAR dans \texttt{\$WILDFLY\_HOME/standalone/deployments/}, ou
déployer depuis l'IDE --- \textbf{mais pas les deux} : un WAR copié à la main
en plus du déploiement de l'IDE provoque l'erreur \texttt{WFLYSRV0205
Duplicate deployment}. Le script \texttt{tools/clean-wildfly-deployment.sh}
nettoie ce cas.

\subsection{La datasource}

L'application n'ouvre \textbf{jamais} de connexion elle-même : elle demande au
serveur la ressource \texttt{java:/OracleDS}. Cette datasource est déclarée
dans WildFly, avec l'URL, le compte et le mot de passe.

\noindent\begin{tabularx}{\textwidth}{|>{\columncolor{btkpale}\bfseries}L{4cm}|X|}
\hline
Nom JNDI & \texttt{java:/OracleDS}\\
\hline
URL & \texttt{jdbc:oracle:thin:@localhost:1521/FREEPDB1}\\
\hline
Compte & celui déclaré côté serveur\\
\hline
Pilote & \texttt{ojdbc11} --- le plus simple est de déposer le \texttt{.jar} dans \texttt{standalone/deployments/}\\
\hline
\end{tabularx}

\begin{retenir}
\textbf{Pourquoi les identifiants ne sont-ils pas dans le projet ?} Parce que
le même WAR doit pouvoir être déployé en développement, en recette et en
production sans être recompilé. Ce qui change d'un environnement à l'autre
appartient au serveur, pas à l'archive. C'est aussi ce qui évite de versionner
un mot de passe.
\end{retenir}

\subsection{Diagnostiquer une erreur de connexion}

L'erreur la plus probable le jour de la soutenance :

\begin{lstlisting}[style=shell]
jakarta.ejb.EJBException: org.hibernate.exception.GenericJDBCException:
  Unable to acquire JDBC Connection
  [jakarta.resource.ResourceException: IJ000453: Unable to get managed
   connection for java:/OracleDS]
\end{lstlisting}

Elle ne vient \textbf{pas} du code : la datasource est bien trouvée, mais son
pool n'arrive pas à joindre Oracle. Une commande tranche entre les deux
familles de causes :

\begin{lstlisting}[style=shell]
python3 etl/test_connexion_oracle.py
\end{lstlisting}

\noindent\begin{tabularx}{\textwidth}{|>{\columncolor{btkpale}\bfseries}L{4.6cm}|X|}
\hline
Le test échoue & Oracle n'est pas joignable : service arrêté, ou base pluggable restée en \texttt{MOUNTED} après un redémarrage --- \texttt{ALTER PLUGGABLE DATABASE FREEPDB1 OPEN;} puis \texttt{SAVE STATE}\\
\hline
Le test réussit & Le problème est la datasource : compte, mot de passe ou nom de service. Corriger par \texttt{jboss-cli}, puis \texttt{reload}\\
\hline
\end{tabularx}

\vspace{0.5em}
La marche à suivre complète est dans
\texttt{tools/DEPANNAGE\_ORACLE\_WILDFLY.md}, et le script de diagnostic dans
\texttt{tools/setup-oracle-datasource.cli}.

\subsection{L'ordre de démarrage, avant la soutenance}

\begin{enumerate}[leftmargin=*,itemsep=1pt]
  \item démarrer Oracle (services Windows) et vérifier que la base pluggable
        est \texttt{READ WRITE} ;
  \item démarrer WildFly ;
  \item ouvrir \texttt{http://localhost:8080/gestion-agences/} ;
  \item si une page a été ouverte pendant qu'Oracle était absent, vider le
        pool : \texttt{flush-all-connection-in-pool}.
\end{enumerate}
